
%%%%%%%%%%%%%%%%%%%%%%%%%%%%%%%%%%%%%
%%      Referencial teórico    %%
%%%%%%%%%%%%%%%%%%%%%%%%%%%%%%%%%%%%%

% Deve haver um tópico para cada assunto abordado na sua pesquisa.

%\chapter{Revisão da Literatura}
\chapter{Questões}
\label{cha:Metodologia}
% remove nome do capítulo do cabeçalho

\section{Questão 1}

\begin{figure}[H]
  \centering
  
\includegraphics[width=1\linewidth]{FIGURAS/mandel-10.jpg}
  % Contém Fonte: 
  %\caption{Questão 01.}
%\source{\cite{xavier2006processamento}}
 % \label{fig:fig01}
\end{figure}
\newpage
\section{Questão 2}

\begin{figure}[H]
  \centering
  
\includegraphics[width=1\linewidth]{FIGURAS/mandel-10.jpg}
  % Contém Fonte: 
  %\caption{Questão 01.}
%\source{\cite{xavier2006processamento}}
 % \label{fig:fig01}
\end{figure}
%\section{Ground Penetrating Radar (GPR)}
\label{subsec:groud_penetrating_radar}

O método aqui descrito é fundamentado no envio de um pulso eletromagnético por meio de uma antena transmissora. Esse pulso tem uma duração extremamente curta de apenas alguns nanosegundos ($10^{-9}$ segundos) empregando ondas eletromagnéticas de frequências VHF/UHF (conforme ilustrado na figura \ref{fig:fig02}).

\begin{figure}[H]
  \centering
  \includegraphics[width=1\linewidth]{lib/fig/gpr/fig01.png}
  % Contém Fonte: 
  \caption{O espectro eletromagnético delineia as gamas de frequência utilizadas por várias tecnologias. O GPR opera dentro da faixa de frequência de 10 MHz a 1 GHz, correspondendo a comprimentos de onda que variam de aproximadamente 30 metros a 0.3 metros.}
\source{\cite{xavier2006processamento}}
  \label{fig:fig01}
\end{figure}

Assim como os métodos sísmicos, o método \ac*{gpr} opera com o registro do tempo que leva para um pulso eletromagnético viajar. Esse pulso é emitido a partir da fonte transmissora (TX), penetra no subsolo e encontra estruturas com contrastes de impedância eletromagnética suficientes para causar reflexões. Em seguida, o pulso retorna ao receptor (RX). A figura 2.2 ilustra o funcionamento do equipamento em um cenário com duas camadas de constantes dielétricas, K1 e K2. A figura também mostra o traço que representa a onda direta pelo ar (Aw), a onda direta através da terra (Gw) e o registro do tempo de viagem duplo até a interface.

\begin{figure}[H]
  \centering
  \includegraphics[width=1\linewidth]{lib/fig/gpr/fig02.png}
  % Contém Fonte: 
  \caption{(A) Princípio de funcionamento do método \ac*{gpr} considerando um ambiente
composto por duas camadas. (B) Traço contendo os registros dos tempos de chegada da
onda direta pelo ar (Aw), on da direta pela terra (Gw) e da onda refletida na interface
formada pelas camadas 1 e 2.}
\source{\cite{oliveira2001dois}}
  \label{fig:fig02}
\end{figure}

Quando esse pulso atinge diferentes meios (ou alvos) com distintas propriedades eletromagnéticas, a velocidade de propagação é alterada. Como resultado, parte da energia eletromagnética é refletida e captada por uma antena receptora, enquanto outra parte sofre difração e continua se propagando até encontrar outra superfície de contraste ou ser completamente atenuada pelo meio.

O "tempo duplo" é calculado, transmitido, registrado e armazenado por uma unidade de controle. Esse tempo duplo refere-se ao intervalo decorrido entre a emissão do pulso eletromagnético pela antena transmissora até o momento em que seu retorno é registrado pela antena receptora, após atingir o alvo e retornar à superfície. Além disso, a unidade de controle também registra a intensidade correspondente do sinal refletido.


\begin{figure}[H]
  \centering
  \includegraphics[width=1\linewidth]{lib/fig/gpr/fig03.png}
  % Contém Fonte: 
  \caption{Esquema de utilização do GPR com transmissão e recebimento de sinais.}
\source{\cite{sampaio2021investigaccao}}
  \label{fig:fig03}
\end{figure}

Dessa forma, é possível obter imagens que revelam informações sobre as condições geomecânicas do terreno ou superfície em estudo, como a densidade de fraturas existentes e seu grau de fraturamento. Além disso, é capaz de determinar o topo do maciço rochoso e identificar descontinuidades, espessura do pacote de sedimentos, textura e nível d'água \cite{davis1989ground}.
Na figura \ref{fig:fig03} mostra o funcionamento do \ac*{gpr}. \cite{annan1992ground}

A validação dos resultados obtidos com o radar evidencia o poder e utilidade desse método como uma ferramenta auxiliar na interpolação e extrapolação das informações obtidas por meio de métodos de perfuração convencionais. \cite{davis1989ground}

\subsection{Janela de GPR}
\label{subsubsec:janela_gpr}

A faixa de operação do \ac*{gpr} consiste em um intervalo de frequência no qual a atenuação experimentada pelo sinal de radar permanece constante para todos os materiais presentes no substrato geológico em consideração. É precisamente dentro dessa faixa que o fenômeno de propagação do pulso eletromagnético é governado pela equação de onda. A frequência associada ao limite inferior dessa faixa deve ser suficientemente alta para evitar a difusão da onda. Por outro lado, o limite superior não pode exceder 1 GHz, uma vez que isso levaria ao fenômeno de polarização das moléculas de água presentes nas rochas, resultando em uma significativa atenuação do sinal \cite{annan1992ground}.

\begin{figure}[H]
  \centering
  \includegraphics[width=1\linewidth]{lib/fig/gpr/fig04.png}
  % Contém Fonte: 
  \caption{A relação entre a velocidade e a frequência em diferentes condutividades.}
\source{\cite{annan1992ground}}
  \label{fig:fig04}
\end{figure}

As velocidades dos sinais de radar em materiais geológicos de baixa perda, adequados para sondagens por radar, estão relacionadas com a parte real da constante dielétrica pela fórmula:
\[V = \frac{c}{(K')^{1/2}}\],
onde \(c = 3 \times 10^8 \, \text{m/s}\), a velocidade de propagação de ondas eletromagnéticas no vácuo.

\begin{table}[H]
    \centering
    \caption{Constante dielétrica típica, condutividade elétrica, velocidade
e atenuação observadas em materiais geológicos comuns a 100
MHz.}

    \begin{tabular}{@{}ccccc@{}}
        \toprule
        \textbf{Material}  & \textbf{Constante
Dielétrica} & \textbf{Condutividade} & \textbf{Velocidade
}      & \textbf{Atenuação}      \\
        \textbf{}  & \textbf{K} & \textbf{$\sigma$(mS/m)} & \textbf{$\nu$(m/ns)
}      & \textbf{$\alpha$(dB/m)}      \\ \midrule
        Ar & 1 & 0 & 0.3 & 0 \\
        Água destilada & 80 & 0.01 & 0.033 & 2 x $10^{-3}$ \\
        Água fresca & 80 & 0.5 & 0.033 & 0.1 \\
        Água do mar & 80 & 3 x $10^{4}$ & 0.01 & 103 \\
        Areia fresca & 3-5 & 0.01 & 0.15 & 0.01 \\
        Areia saturada & 20-30 & 0.1-1 & 0.06 & 0.03-0.3 \\
        Calcário & 4-8 & 0.5-2 & 0.12 & 0.4-1 \\
        Folhelho & 5-15 & 1-100 & 0.09 & 1-100 \\
        Siltito & 5-30 & 1-100 & 0.07 & 1-100 \\
        Argila & 5-40 & 2-1000 & 0.06 & 1-300 \\
        Granito & 4-6 & 0.01-1 & 0.13 & 0.01-1 \\
        Sal & 5-6 & 0.01-1 & 0.13 & 0.01-1 \\
        Gelo & 3-4 & 0.01 & 0.16 & 0.01 \\
        \multicolumn{1}{l}{} & \multicolumn{1}{l}{}         & \multicolumn{1}{l}{} & \multicolumn{1}{l}{} \\ \bottomrule
    \end{tabular}
    \source{\cite{daniels2004ground}}
    \label{tab:tab1}
\end{table}

\subsection{GPRPy}
\label{subsubsec:gprpy}

Muitas vezes, as investigações envolvendo \ac{gpr} não exigem um processamento avançado para obter resultados úteis. No entanto, a maioria dos estudos de pesquisa em \ac{gpr} utiliza software proprietário fechado fornecido pelos fabricantes de sistemas \ac{gpr} ou pacotes independentes \citeonline{jol2008ground}. Essa dependência em software proprietário apresenta várias desvantagens significativas. Em primeiro lugar, o software oferece pouca flexibilidade, impedindo os usuários de adaptá-lo completamente às suas necessidades específicas. Em segundo lugar, a capacidade de reproduzir pesquisas é limitada, pois para replicar os passos de processamento mencionados em um artigo, os pesquisadores precisam ter acesso ao mesmo software proprietário. Essa limitação se torna mais crítica à medida que editores e agências de financiamento pressionam por pesquisas transparentes. Em terceiro lugar, a disponibilidade de cursos que ensinam o processamento de \ac{gpr} também é prejudicada devido à dependência do software proprietário.

O pacote de software \ac{gprpy} consiste em dois interfaces gráficas de usuário (GUIs): uma para o processamento de dados de perfil figura \ref{fig:fig07} e outra para dados de ponto médio comum (CMP) ou reflexão e refração de amplo ângulo (WARR) figura \ref{fig:fig07}. O pacote também inclui uma série de funções para processamento mais avançado, como a geração de cubos de dados 3D por interpolação de perfis, interpolação de superfícies entre pontos selecionados e fusão de perfis. O código de computador da GUI é separado do código de computador que implementa o processamento e a visualização de dados. Essa separação permite que os usuários escrevam código Python que acesse diretamente todas as capacidades do \ac{gprpy} sem precisar usar a GUI. O \ac{gprpy} aproveita essa separação por meio de seus scripts gerados automaticamente.

\begin{figure}[H]
  \centering
  \includegraphics[width=1\linewidth]{lib/fig/gpr/fig07.png}
  % Contém Fonte: 
  \caption{Interface Gráfica do Usuário (GUI) para o processamento de dados de perfil de ponto médio comum (CMP) no \ac{gpr}.}
\source{\cite{plattner2020gprpy}}
  \label{fig:fig07}
\end{figure}

\begin{figure}[H]
  \centering
  \includegraphics[width=1\linewidth]{lib/fig/gpr/fig08.png}
  % Contém Fonte: 
  \caption{Interface Gráfica do Usuário (GUI) para a análise de velocidade a partir de dados CMP ou WARR, com representações de empilhamento de amplitudes para padrões de tempo de chegada em forma de hipérbole (a) e em forma linear (b), bem como os próprios dados CMP ou WARR (c).}
\source{\cite{plattner2020gprpy}}
  \label{fig:fig08}
\end{figure}

\section{ANÁLISE E INTERPRETAÇÃO DE DADOS}
\label{subsubsec:analiseeinterpretação}

O processamento de dados do \ac*{gpr} pode seguir dois caminhos principais. O primeiro envolve a apresentação dos dados em forma de seções, planos ou volumes para indicar a localização de alvos anômalos, comum a todos os métodos geofísicos. O segundo caminho consiste em extrair variáveis das propriedades das ondas, como velocidade, atenuação ou impedância, e convertê-las em informações específicas para a aplicação em questão.

O processo típico de processamento de dados do \ac*{gpr} inclui etapas de edição de dados, processamento básico, processamento avançado e visualização/interpretação. O processamento é uma atividade iterativa, e os dados podem passar várias vezes por esse ciclo. Em grandes conjuntos de dados, pode ser aplicado um processamento em lote com controle interativo limitado, após testes iniciais em amostras selecionadas.

\begin{figure}[H]
  \centering
  \includegraphics[width=1\linewidth]{lib/fig/gpr/fig05.png}
  % Contém Fonte: 
  \caption{Visão geral do fluxo de processamento de dados do radar de penetração no solo (GPR).}
\source{\cite{jol2008ground}}
  \label{fig:fig05}
\end{figure}

\begin{itemize}
    \item Instrumento: Equipamentos ou dispositivos usados para coletar dados ou medições;
    \item Impressão: Gravação de dados que são exibidos ou impressos de alguma forma;
    \item Exibição/Impressão: Exibir ou imprimir dados, para fins de visualização ou análise;
    \item Visualização: Ferramentas ou software usados para criar representações visuais de dados, tornando mais fácil compreendê-los e analisá-los;
    \item Editar: Processo de edição ou modificação de dados, o que pode ser necessário para limpar ou pré-processar os dados;
    \item Processo Básico: Etapas fundamentais de processamento de dados, como filtragem, redução de ruído ou cálculos simples;
    \item Processo Avançado: Técnicas de processamento de dados mais complexas que vão além das operações básicas. Pode incluir tarefas como extração de características, reconhecimento de padrões ou análise estatística;
    \item Processamento de Imagens: Imagem processada, após tratamento básico e avançado, gerando imagem imagem geológica ou não, para uma melhor interpretação dos dados;
    
\end{itemize}

\begin{algorithm}[H]
  \caption{Pseudocódigo da Coletar e Processamento de dados brutos do \ac{gpr}.}
  \label{alg:pgpr}
  \SetAlgoLined
  \begin{algorithmic}[1]
    \STATE Coletar dados brutos do \ac*{gpr}
    \STATE Pré-processamento dos dados:
    \STATE Inicializar uma imagem vazia para armazenar a imagem otimizada.
    \FOR{$i=1$ até $i=4$ (i)}
    \STATE Inicializar uma imagem vazia para armazenar a imagem otimizada.
    \ENDFOR
    \STATE Pré-processamento:
    \STATE Realizar filtragem de imagem adicional.
    \STATE Aplicar técnicas de realce de recursos.
    \STATE Gerar visualizações finais da imagem otimizada.
    \STATE Atualize T se houver uma solução melhor
    \STATE Exibir a imagem otimizada para análise e interpretação.
    \STATE Salvar a imagem otimizada, se necessário.
    \STATE \Return{Imagem processada.}
  \end{algorithmic}
\end{algorithm}

%1. Coletar dados brutos do GPR.

%2. Pré-processamento dos dados:
%   a. Aplicar correções de ganho e calibração.
 %  b. Remover ruído indesejado.
  % c. Realizar correções de distorção.

%3. Inicializar uma imagem vazia para armazenar a imagem otimizada.

%4. Para cada ponto na grade da imagem:
 %  a. Definir uma região de interesse (ROI) em torno do ponto.
  % b. Calcular métricas de qualidade da imagem dentro da ROI, como contraste e nitidez.
   %c. Ajustar os parâmetros de processamento, como filtro e equalização, com base nas métricas de qualidade.
   %d. Aplicar os parâmetros de processamento otimizados à ROI.
   %e. Atualizar a imagem otimizada com os resultados do processamento.

%5. Repetir o passo 4 até que todos os pontos da imagem tenham sido processados.

%6. Pós-processamento opcional:
 ¨%  a. Realizar filtragem de imagem adicional.
 %  b. Aplicar técnicas de realce de recursos.
  % c. Gerar visualizações finais da imagem otimizada.

%7. Exibir a imagem otimizada para análise e interpretação.

%8. Salvar a imagem otimizada, se necessário.

%9. Fim do algoritmo.

Métodos avançados de processamento de dados do \ac*{gpr}, como filtragem espacial e temporal, silenciamento seletivo, filtragem de inclinação, deconvolução e análise de semelhança de velocidade, resultam em dados substancialmente diferentes das informações brutas. Além disso, operações específicas para o \ac*{gpr}, como remoção de fundo e mistura de antenas de múltiplas frequências, são aplicadas.

A evolução do \ac*{gpr} é documentada nos Anais de Conferências de \ac*{gpr} e fornece informações valiosas sobre o desenvolvimento do processamento de dados do \ac*{gpr}. A discussão aqui se concentra em questões essenciais do \ac*{gpr}, que diferem significativamente do processamento sísmico.

\subsection{Dewow}
\label{subsubsec:dewow}

Processo que é usado para suprimir um efeito chamado "wow" que ocorre nos dados do Radar de Penetração no Solo \ac{gpr} devido à proximidade do receptor com o transmissor. O "wow" é causado pela energia de baixa frequência nos campos próximos ao transmissor, que faz com que o nível base do sinal recebido se curve para cima ou para baixo. O "dewow" envolve a aplicação de um filtro temporal de alto amortecimento para eliminar esse efeito, e pode ser realizado em sistemas analógicos mais antigos ou digitalmente em sistemas modernos.

\subsection{Ganho de Tempo}
\label{subsubsec:ganhodetempo}

O "ganho de tempo", que é usado para equalizar as amplitudes dos sinais de radar de \ac*{gpr}, uma vez que esses sinais se atenuam rapidamente à medida que se propagam no solo. O ganho de tempo envolve a aplicação de uma função de ganho dependente do tempo para compensar essa rápida queda nos sinais de radar de maiores profundidades. A atenuação no solo pode variar significativamente, portanto, o ganho de tempo é aplicado para acomodar essas variações e permitir a exibição adequada dos dados de \ac*{gpr} versus tempo. O ganho de tempo pode ser aplicado de várias maneiras e deve ser escolhido com base em um modelo físico para minimizar artefatos no processo.

\subsection{Deconvolução}
\label{subsubsec:deconvolução}

Processo que visa maximizar a largura de banda e reduzir a dispersão de pulsos para aumentar a resolução em dados de \ac*{gpr}. No entanto, geralmente, a deconvolução de dados de \ac*{gpr} não oferece muitos benefícios significativos, pois o pulso padrão do \ac*{gpr} já é o mais curto e comprimido possível para as condições dadas de largura de banda e relação sinal-ruído. A deconvolução pode ser útil em casos de reverberação extrínseca ou reverberação do sistema.

Um tópico relacionado é a filtragem inversa Q, que tenta compensar a rápida atenuação das frequências mais altas em dados de \ac*{gpr} à medida que a profundidade aumenta. Isso é importante para manter a resolução em maiores profundidades. \citeonline{irving2003removal} fornecem uma discussão detalhada sobre a filtragem inversa Q.

\subsection{Migração}
\label{subsubsec:migração}


A migração é uma técnica de deconvolução espacial que visa remover a direcionalidade da fonte e do receptor dos dados de reflexão. Seu objetivo é reconstruir uma representação precisa e geométrica da distribuição de refletividade do subsolo, mas para isso, é necessário ter conhecimento da estrutura de velocidade. A migração é frequentemente um processo interativo, onde a velocidade de fundo é ajustada iterativamente para otimizar a imagem resultante conforme figura \ref{fig:fig06} . Existem várias formas de migração possíveis, e seu uso adequado no \ac*{gpr} visa compensar as características direcionais das antenas. Embora muitos princípios de migração sísmica se apliquem ao \ac*{gpr}, esse método também lida com o desafio único da topografia em relação à profundidade da exploração. A consideração da topografia é discutida por \citeonline{lehmann2000topographic}.

\begin{figure}[H]
  \centering
  \includegraphics[width=1\linewidth]{lib/fig/gpr/fig06.png}
  % Contém Fonte: 
  \caption{Imagens obtidas com um sistema de 50 MHz que atravessam dois túneis rodoviários, antes e após a migração.}
\source{\cite{jol2008ground}}
  \label{fig:fig06}
\end{figure}

\subsection{Correção topográfica}
\label{subsubsec:topografica}

Devido à pequena profundidade de exploração do \ac*{gpr}, a compensação da topografia frequentemente é importante. Para variações superficiais menores, o deslocamento temporal das trajetórias de dados pode compensar em grande parte as variações topográficas. O exemplo de perfil de reflexão na figura \ref{fig:fig10}  mostra a estratigrafia e uma reflexão da superfície da água. O perfil começa em uma área plana e sobe uma encosta. Os dados foram compensados por deslocamento temporal em relação à elevação ao longo da linha do perfil, o que torna a reflexão da superfície da água essencialmente plana.

\begin{figure}[H]
  \centering
  \includegraphics[width=1\linewidth]{lib/fig/gpr/fig10.png}
  % Contém Fonte: 
  \caption{Ilustração usando o deslocamento temporal dos traços para compensar variações topográficas. Observe que o nível freático é essencialmente plano na seção de radar de penetração no solo \ac{gpr}.}
\source{\cite{jol2008ground}}
  \label{fig:fig10}
\end{figure}
%
\section{Método da eletrorresistividade}
\label{subsec:metodo_da_eletrorresistividade}

A \ac{er} é um método geofísico utilizado para determinar a resistividade elétrica dos diferentes tipos de rochas. \cite{braga_metodos_2006}

% Resistividadee
A resistividade elétrica de um material \ref{fig:fig01} é definida pela sua resistência em ohms. Para um cilindro condutor de resistência $\delta R$, comprimento $\delta L$ e área de seção transversal $\delta A$ (Figura \ref{fig:ro}), a resistividade $\rho$ é dada por:
\begin{align}
  \label{eq:resistivity}
  \rho = \frac{\delta{R}\delta{A}}{\delta{L}}
\end{align}

\begin{figure}[H]
  \centering
  \includegraphics[width=1\linewidth]{lib/fig/resistividade/ro.png}
  % Contém Fonte: 
  \caption{Parâmetros usados na definição de resistividade}
  \source{\cite{Kearey_introduction_2002}}
  \label{fig:ro}
\end{figure}

Diferentes tipos de minerais podem conduzir melhor eletricidade do que outros, assim a medida de resistividade pode ser calculada utilizando a equação \ref{eq:resistivity}. Na figura \ref{fig:res_mat} apresenta os valores dos principais tipos de rochas com suas resistividades.

\begin{figure}[H]
  \centering
  \includegraphics[width=1\linewidth]{lib/fig/resistividade/resistividade.png}
  \caption{Intervalo aproximado de valores de resistividade de alguns materiais}
  \source{\cite{de2016geofisica}}
  \label{fig:res_mat}
\end{figure}

% Resistividade aparente

% SEV
A \ac{sev} ou \ac{ves} é utilizado um equipamento para injetar corrente elétrica no solo, a partir de diferentes espaçamentos entre os eletrodos de corrente (A e B) dispostos na superfície do terreno e outros dois eletrodos de potencial (M e N), calcula-se a diferença de potencial entre eles, segundo o arranjo Schlumberger. Na Figura \ref{fig:sev} contém medidas de resistividade ($\rho_n$) e espessuras das camadas ($E_n$) em uma \ac{sev} fictícia. \cite{de2016geofisica}

\begin{figure}[H]
  \centering
  \includegraphics[width=1\linewidth]{lib/fig/resistividade/sev_dados1.jpg}
  % Contém Fonte: 
  \caption{Sondagem Elétrica Vertical}
  \source{Elaboração própria a partir de \cite{de2016geofisica}}
  \label{fig:sev}
\end{figure}


Quando sua investigação é programada para uma profundidade fixa, as medidas resultam num perfil geoelétrico em uma dimensão (1D) e seus resultados obtidos podem ser expressos por meio de uma curva de resistividade aparente ($\rho_a$) em Ohm metro ($\Omega m$), Figura \ref{fig:curva}.

\begin{figure}[H]
  \centering
  \includegraphics[width=1\linewidth]{lib/fig/resistividade/curva.png}
  \caption{Curva de resistividades aparentes.}
  % Contém Fonte: 
  
  \label{fig:curva}
\end{figure}
%\section{Inversão de dados em eletrorresistividade}
\label{subsec:inversao_de_dados_em_er}

A Teoria da inversão de dados é utilizada em eletrorresistividade para estimar modelos geológicos que justifiquem os valores das medidas físicas adquiridas durante levantamentos efetuados na superfície. Os dados coletados no campo podem ser expressas através de um vetor de dados $d = {[~d_1,~d_2,~d_3~...~d_n~]}^{T}$, onde $T$ significa transposta. Da mesma forma, os parâmetros do modelo podem ser representados como os elementos de um vetor modelo $m = [~\rho_1,~\rho_2,~\rho_3~...~\rho_n,~E_1,~E_2,~E_3~...~E_n-1]^{T}$. O parâmetro do modelo ($2n - 1$) é a última camada, sendo o semi-espaço. Os parâmetros $\rho_n$ e $E_n$ são os valores de resistividade e espessura, respectivamente.

Esses vetores são relacionados através de um operador de modelagem $g$, conforme equação abaixo \cite{koefoed1979resistivity}:
\begin{align}
    \label{eq:dgm}
    d=g(m)
\end{align}

No processo de inversão, estimamos um modelo $m_{est}$ que minimiza a função objetivo ou função \textit{fitness}, a qual é uma medida de discrepância entre os dados observados e dados calculados. Conforme a equação abaixo:
\begin{equation}
    \label{eq:fit}
    F=\frac{1}{N}|| d^{obs}-d^{cal}||^2
\end{equation}


Os dados observados ($d_{obs}$) e os dados calculados ($d_{calc}$) bem como o número total de dados ($N$) são calculados, para dados sintéticos e reais, usando a equação \ref{eq:fit}. \cite{Benjumea_1d_2021}

Na inversão deve ser utilizado uma curva de campo para decidir o número de camadas e resistividades aparente e definir um limite de pesquisa de valores de resistividade e espessura da camada. \citeonline{Sen1993} utilizam uma abordagem que serve como direcionamento para o processo de inversão desse trabalho.
%\section{Algoritmo de otimização do gafanhoto}
\label{subsec:algoritmo-do-gafanhoto}
Proposto por \citeonline{Saremi2017}, o \acf{aog} ou \ac{goa} é um algoritmo que simula o comportamento do gafanhoto em enxame na busca por comida na natureza. Os gafanhotos são insetos conhecidos como uma das pragas mais perigosas que dizimam plantações e prejudicam a agricultura em algumas regiões do mundo. Seu ciclo de vida inclui duas fases: a fase de ninfa (figura \ref{fig:ninfa}) é caracterizada pelos seus movimentos lentos e pequenos passos (explotação), já na fase adulta (figura \ref{fig:adulto}), seus movimentos são bruscos e de longo alcance (exploração). Esses movimentos constituem a base para o algoritmo para o \ac{aog}.

\begin{figure}[H]
  \centering
  \includegraphics[width=1\linewidth]{lib/fig/gafanhoto/ninfa.jpg}
  \caption{Gafanhotos na fase de ninfa}
  \source{\cite{noauthor_desert_2016}}
  \label{fig:ninfa}
\end{figure}

\begin{figure}[H]
  \centering
  \includegraphics[width=1\linewidth]{lib/fig/gafanhoto/adulto.png}
  % Contém Fonte: 
  \caption{Gafanhotos na fase adulta}
  \source{\cite{kaufman_why_2020}}
  \label{fig:adulto}
\end{figure}

A equação \ref{eq:pi} simula o comportamento de enxame dos gafanhotos:

\begin{equation}
  \label{eq:pi}
  P_i = S_i + G_i + A_i
\end{equation}
onde $P_i$ indica a posição do i-ésimo gafanhoto, $S_i$ representa a interação social entre a solução e os outros gafanhotos, $G_i$ é a força da gravidade sobre o i-ésimo gafanhoto e $A_i$ representa a direção do vento.
Para produzir um comportamento aleatório dos gafanhotos, a equação \ref{eq:pi} pode ser escrita como:

\begin{equation}
    \label{eq:pi_r}
    P_i = r_1 S_i + r_2 G_i + r_3 A_i
  \end{equation}
  onde $r_1$, $r_2 $ e $r_3$ são números aleatórios entre $[0,1]$.

\subsection{Interação Social}
\label{subsubsec:interacao_social}

  A interação social entre os gafanhotos $S_i$ é definida como:
  \begin{equation}
    \label{eq:si}
    S_i = \sum_{\substack{j=1 \\i \neq j }}^{N} ~s(d_{ij})~ \widehat{d_{ij}}
  \end{equation}

    \begin{equation}
    \label{eq:sr}
    s = f~ exp^{-r/l} - exp^{-r}
  \end{equation}
  
  
  
  % \begin{equation}
  %   \label{eq:dij}
  %   d_{ij} = | P_j - P_i |
  % \end{equation}
  
  
  % \begin{equation}
  %   \label{eq:dijv}
  %   \widehat{d_{ij}} =\frac{P_j-P_i}{d_{ij}}
  % \end{equation}
  

  
  Onde $N$ na equação \ref{eq:si} refere-se ao número de gafanhotos, onde $ d_{ij} = | P_j - P_i |$ representa a distância entre o i-ésimo e o j-ésimo gafanhoto, $\widehat{d_{ij}} =\frac{P_j-P_i}{d_{ij}}$ representa o vetor unitário.
  
A equação \ref{eq:sr} representa a força de repulsão e atração entre gafanhotos, onde $l$ é a escala de comprimento atrativo e $f$ é a intensidade da atração. Quando a distância entre dois gafanhotos  é entre [0, 2.079] ocorre repulsão, e quando a distância entre dois gafanhotos é exatamente 2.079, não há nem repulsão, nem atração (sem força), esta área é chamada de zona de conforto. Quando a distância excede 2.079, a força de atração aumenta, depois diminui progressivamente até atingir 4. O valor da função $s$ é próximo de zero e a distância é maior que 10. Portanto, não podemos usar esta função para forças fortes entre gafanhotos com grandes distâncias entre eles.

A figura abaixo mostra a interação entre gafanhotos em relação à zona de conforto.
  
  \begin{figure}[H]
    \centering
    \includegraphics[width=1\linewidth]{lib/fig/gafanhoto/atracao_repul.png}
    % Contém Fonte: 
    \caption{interação entre os gafanhotos. Gafanhotos de cor vermelha estão sendo atraídos para a zona de conforto [0,2079], gafanhotos de cor verde estão na zona de conforto [2.079] e gafanhotos de cor azul estão em repulsão para a zona de conforto [2.079,4].}
    \source{O Autor}
    \label{fig:atra_repu}
  \end{figure}
  

\subsection{Força da Gravidade}
\label{subsubsec:forca_gravidade}

A força da gravidade $G_i$ é dada pela equação:

\begin{equation}
  \label{eq:gi}
  G_i = -g \hat{e}_g
\end{equation}
onde $g$ representa a constante gravitacional e $\hat{e}_g $ representa um vetor unitário em direção ao centro da Terra.


\subsection{Direção do Vento}
\label{subsubsec:direcao_vento}

A direção do vento $A_i$ é dada pela equação:

\begin{equation}
  \label{eq:ai}
  A_i = u \hat{e}_w
\end{equation}
onde $u$ representa uma constante de deriva e $\hat{e}_w$ é um vetor unitário na direção do vento.
\subsection{Posição do Gafanhoto}
\label{subsubsec:posicao_gafanhoto}


Quando substituímos os valores na equação \ref{eq:si}, temos:

\begin{equation}
  {P_i} = \sum_{\substack{j=1 j \neq i}}^{N} s~(d_{ij})\widehat{d_{ij}} - g\hat{e}_g + u\hat{e}_w = \sum_{\substack{j=1 j \neq i}}^{N} s~(|P_j-P_i|) \frac{P_j-P_i}{d_{ij}}- g\hat{e}_g + u\hat{e}_w
  \label{eq:pi_mod}
\end{equation}

No entanto, \citeonline{Saremi2017} mostram que a equação \ref{eq:pi_mod} não pode ser usada diretamente para resolver problemas de otimização, pois os gafanhotos chegam rapidamente à zona de conforto e o sistema de enxame não converge para sua solução (ótimo global). Então a equação é alterada para:

\begin{equation}
  \label{eq:pid}
  P_i^d = ~c\left (\sum_{\substack{j=1 \\j \neq i}}^{N} ~c~\frac{ub_d-lb_d}{2} s\big (|P_j^d-P_i^d| \big ) \frac{P_j-P_i}{d_{ij}}\right ) + \widehat{T}_d
\end{equation}
onde $ub_d$ e $lb_d$ representam os limites superior e inferior na d-ésima dimensão, respectivamente. $\hat{T}_d$ é a melhor solução encontrada até o momento no espaço de dimensão $d$ e $G_i = 0$.

O parâmetro $c_1$ é usado para reduzir os movimentos dos gafanhotos em torno da solução, proporcionando um bom equilíbrio entre intensificação e diversificação. O parâmetro $c_2$ é usado para reduzir a zona de repulsão, zona de atração e zona de conforto entre gafanhotos proporcionalmente ao número de iterações. Para equilibrar as fases de exploração e explotação usando a abordagem do gafanhoto, o coeficiente $c$ diminui conforme o número de iterações.

\begin{equation}
  \label{eq:c}
  c = c_{max} - iter~\frac{c_{max} - c_{min}}{Max_{iter}}
\end{equation}
onde $c_{max}$ e $c_{min}$ representam o máximo e o mínimo valores de $c$, respectivamente, $t$ é a iteração atual e $Max_{iter}$ é o número máximo de iterações.

\subsection{Pseudocódigo e Complexidade}
\label{subsubsec:pseudocodigo_comlexidade}

O pseudocódigo do \ac{aog} é mostrado no algoritmo \ref{alg:paog} e o código-fonte, em várias linguagens de programação, pode ser encontrado na página \href{https://seyedalimirjalili.com/goa}{https://seyedalimirjalili.com\\/goa}.

\begin{algorithm}[H]
  \caption{Pseudocódigo do Algoritmo de
    Otimização do Gafanhoto (AOG)}
  \label{alg:paog}
  \SetAlgoLined
  \begin{algorithmic}[1]
    \STATE Gere a população inicial de Gafanhotos aleatoriamente $P_i(i=1,2,...,n)$
    \STATE Inicialize $c_{min}$, $c_{max}$ e número máximo de iterações $Max_{iter}$
    \STATE Avalie o fitness $F(P_i)$ de cada gafanhoto $P_i$
    \STATE T = a melhor solução
    \WHILE{$(iter<Max_{iter})$}
    \STATE Atualize $c_1$ e $c_2 $ usando a equação \ref{eq:c}
    \FOR{$i=1$ até $N$ (todos N gafanhotos na população)}
    \STATE Normalize a distância entre gafanhotos no intervalo [1,4]
    \STATE Atualize a posição do gafanhoto atual usando a equação \ref{eq:pid}
    \STATE Traga o gafanhoto atual de volta se ele for fora dos limites
    \ENDFOR
    \STATE Atualize T se houver uma solução melhor
    \STATE $iter = iter + 1$
    \ENDWHILE
    \STATE \Return{a melhor solução T}
  \end{algorithmic}
\end{algorithm}

O \ac{aog} começa gerando a população de gafanhotos e inicializando $c_{min}$, $c_{max}$, e número máximo de interações ($Max_{iter}$). Em seguida, avalia-se cada solução calculando seu valor usando a função \textit{fitness}. Após avaliar cada solução na população, atribuí-se a melhor solução de acordo com seu valor.
Depois, atualiza-se o parâmetro do coeficiente para encolher as zonas de atração, repulsão e conforto (Equação  \ref{eq:c}).

A função $s$ da Equação \ref{eq:sr}, divide o espaço de busca em zonas de
repulsão, conforto e atração. Quando essa função falha em aplicar as forças entre os gafanhotos, mapeia-se a distância no intervalo {[1,4]}. 

Em seguida, atualiza-se cada solução na população com base na distância entre o gafanhoto (a solução) e o outro gafanhoto (a outra solução). Também atualiza-se a melhor solução na população e o parâmetro do coeficiente decrescente usando a Equação \ref{eq:pid}.

No caso de um gafanhoto violar os limites após atualiza-se a solução, traz o gafanhoto de volta ao seu domínio. Em seguida, repete-se os três passos anteriores para cada solução na população.

Depois, atualiza e avalia cada solução para atribuir a melhor na população. Repete as operações gerais até atingir o número máximo de iterações para retornar a melhor solução global.

A complexidade de tempo do \ac{aog} está diretamente relacionada ao número de agentes($n$) e número de interações($n \times Max_{iter}$). Assim, $O(n \times Max_{iter})$ representa a complexidade de tempo desse algoritmo. \cite{GOA_web}
%% \section{Modelagem em eletrorresistividade 1D}
% \label{subsec:modelagem-em-eletrorresistividade 1D}
